% titel: Blatt 7, A3: Staumauer
% skript: Physik I, Kapitel 8.1.1--8.1.2, Seiten 129--133
% stand: Übungsblatt 7, Aufgabe A3

Wasser drückt auf die Mauer, die das Becken hält. Gesucht wird, wie gross
diese Kraft ist und wo sie angreift. Es genügt die Hydrostatik einer ruhenden
Flüssigkeit; Bewegung oder Strömung spielen keine Rolle.

\begin{werkzeugkasten}[Druck in einer ruhenden Flüssigkeit]
Druck ist Kraft senkrecht auf die Fläche. In ruhender Flüssigkeit konstanter
Dichte $\rho$ wächst er mit der Tiefe $h$ unter der freien Oberfläche linear an:
\begin{equation}
  p(h)=p_0+\rho gh.
\end{equation}
$p_0$ ist der Druck an der Oberfläche, $\rho gh$ der Schweredruck. Der Druck
wirkt in alle Richtungen, auch auf Seitenwände, und wächst mit der Tiefe.
Deshalb ist er an der Mauer nicht konstant, sondern unten am grössten.

\textit{Skript:} Kap. 8.1.1--8.1.2, S. 129--132, Gl. (8.4).
\end{werkzeugkasten}

\begin{werkzeugkasten}[Vom Druck zur Kraft auf die Mauer]
Für die Belastung zählt der Druckunterschied zwischen den beiden Seiten. Die
Kraft auf ein Flächenstück ist Druck mal Fläche; da der Druck über die Höhe
variiert, muss man die Beiträge über die benetzte Fläche aufsummieren
(integrieren), statt pauschal ``Druck mal Gesamtfläche'' zu rechnen. Der
Druck selbst ist von der Gefässform unabhängig; die Kraft auf die Wand hängt
von deren Gesamtfläche und der Verteilung ab.

\textit{Skript:} Kap. 8.1.2, S. 131--133.
\end{werkzeugkasten}

\begin{achtung}[Darauf achten]
\begin{itemize}
  \item Tiefe unter dem Wasserspiegel und Höhe über dem Boden sauber trennen;
        $p$ hängt nur von der senkrechten Tiefe ab.
  \item Absolute Höhe und Wasserspiegelhöhe passend einsetzen; an der Mauer
        beginnt die Tiefe oben bei null.
\end{itemize}
\end{achtung}
